% GENERATED FILE. Edit canonical lessons, then run npm run generate:books.
% canonical-source-hash: fnv1a64:dc75318519f60390
% canonical-lessons: JA-C217-kanja, JA-C217-shinsatsu, JA-C217-chuusha, JA-C217-haisha, JA-C217-mushiba

\chapter{At the Doctor's}
\label{ch:ja-at-the-doctor-s}
\hlchaptermodality{217}

\begin{chapteropening}
\textbf{By the end of this chapter:} \emph{I can answer a doctor's simple questions.}

Complete the last lesson of chapter 217: I can answer a doctor's simple questions.
\end{chapteropening}

\section[\texorpdfstring{\ja{かんじゃ} (kanja)}{かんじゃ (kanja)}]{\ja{かんじゃ} (kanja) — a patient}

\label{lesson:JA-C217-kanja}



\begin{quote}
Before the new one: say the Japanese for to get better, then the Japanese for to shiver.
\end{quote}

\subsection*{You'll want to know: \ja{かんじゃ}}
\textbf{\ja{かんじゃ}} — \emph{kanja} — \textquotedblleft{}a patient\textquotedblright{}.

\begin{quote}
\textbf{\ja{びょういんに} \ja{かんじゃは} \ja{おおいですか。}} — \emph{byōin ni kanja wa ōi desu ka} — \textquotedblleft{}Are there many patients at the hospital?\textquotedblright{}
\end{quote}

\begin{cousinweb}
\textbf{\ja{かんじゃ}} (\emph{kanja}) is Sino-Japanese (kango): two Chinese characters read with their on'yomi, readings borrowed from Chinese. They are \emph{kan}, \textquotedblleft{}to suffer, illness\textquotedblright{} (Middle Chinese \emph{hwaenH}) and \emph{ja}, \textquotedblleft{}person\textquotedblright{} (\emph{tsyaeX}) (KANJIDIC2; Middle Chinese in Baxter's transcription). The second is voiced from \emph{sha} to \emph{ja}.
\end{cousinweb}

\begin{grammarlens}[title={What you've built}]
One more word for everyday life at A2.
\end{grammarlens}

\subsection*{Your turn}
\begin{itemize}
\raggedright
  \item \emph{Say it:} \emph{kanja}
  \item \emph{Say it:} \emph{kanja}, once more
  \item \emph{Say it:} \emph{byōin ni kanja wa ōi desu ka} — Are there many patients at the hospital?
  \item \emph{Recall:} say the Japanese for to get better, then the Japanese for to shiver, then say \emph{kanja} again
\end{itemize}

\begin{tcolorbox}[breakable,colback=teal!4,colframe=teal!35!black,title={Before you move on}]
What does \ja{かんじゃ} mean? (\textquotedblleft{}A patient\textquotedblright{}.) Say it once more.
\end{tcolorbox}
\section[\texorpdfstring{\ja{しんさつ} (shinsatsu)}{しんさつ (shinsatsu)}]{\ja{しんさつ} (shinsatsu) — a consultation; a medical examination}

\label{lesson:JA-C217-shinsatsu}



\begin{quote}
Before the new one: say the Japanese for to shiver, then the Japanese for a patient.
\end{quote}

\subsection*{You'll want to know: \ja{しんさつ}}
\textbf{\ja{しんさつ}} — \emph{shinsatsu} — \textquotedblleft{}a consultation; a medical examination\textquotedblright{}.

\begin{quote}
\textbf{\ja{しんさつは} \ja{ごごですか。}} — \emph{shinsatsu wa gogo desu ka} — \textquotedblleft{}Is the consultation in the afternoon?\textquotedblright{}
\end{quote}

\begin{cousinweb}
\textbf{\ja{しんさつ}} (\emph{shinsatsu}) is Sino-Japanese (kango): two Chinese characters read with their on'yomi, readings borrowed from Chinese. They are \emph{shin}, \textquotedblleft{}to examine\textquotedblright{} (Middle Chinese \emph{tsyinX}) and \emph{satsu}, \textquotedblleft{}to judge, understand\textquotedblright{} (\emph{tsrheat}) (KANJIDIC2; Middle Chinese in Baxter's transcription).
\end{cousinweb}

\begin{grammarlens}[title={What you've built}]
One more word for everyday life at A2.
\end{grammarlens}

\subsection*{Your turn}
\begin{itemize}
\raggedright
  \item \emph{Say it:} \emph{shinsatsu}
  \item \emph{Say it:} \emph{shinsatsu}, once more
  \item \emph{Say it:} \emph{shinsatsu wa gogo desu ka} — Is the consultation in the afternoon?
  \item \emph{Recall:} say the Japanese for to shiver, then the Japanese for a patient, then say \emph{shinsatsu} again
\end{itemize}

\begin{tcolorbox}[breakable,colback=teal!4,colframe=teal!35!black,title={Before you move on}]
What does \ja{しんさつ} mean? (\textquotedblleft{}A consultation; a medical examination\textquotedblright{}.) Say it once more.
\end{tcolorbox}
\section[\texorpdfstring{\ja{ちゅうしゃ} (chūsha)}{ちゅうしゃ (chūsha)}]{\ja{ちゅうしゃ} (chūsha) — an injection}

\label{lesson:JA-C217-chuusha}



\begin{quote}
Before the new one: say the Japanese for a patient, then the Japanese for a consultation.
\end{quote}

\subsection*{You'll want to know: \ja{ちゅうしゃ}}
\textbf{\ja{ちゅうしゃ}} — \emph{chūsha} — \textquotedblleft{}an injection\textquotedblright{}.

\begin{quote}
\textbf{\ja{ちゅうしゃは} \ja{いたいですか。}} — \emph{chūsha wa itai desu ka} — \textquotedblleft{}Does the injection hurt?\textquotedblright{}
\end{quote}

\begin{cousinweb}
\textbf{\ja{ちゅうしゃ}} (\emph{chūsha}) is Sino-Japanese (kango): two Chinese characters read with their on'yomi, readings borrowed from Chinese. They are \emph{chū}, \textquotedblleft{}to pour\textquotedblright{} (Middle Chinese \emph{tsyuH}) and \emph{sha}, \textquotedblleft{}to shoot\textquotedblright{} (\emph{zyaeH}) (KANJIDIC2; Middle Chinese in Baxter's transcription).
\end{cousinweb}

\begin{grammarlens}[title={What you've built}]
One more word for everyday life at A2.
\end{grammarlens}

\subsection*{Your turn}
\begin{itemize}
\raggedright
  \item \emph{Say it:} \emph{chūsha}
  \item \emph{Say it:} \emph{chūsha}, once more
  \item \emph{Say it:} \emph{chūsha wa itai desu ka} — Does the injection hurt?
  \item \emph{Recall:} say the Japanese for a patient, then the Japanese for a consultation, then say \emph{chūsha} again
\end{itemize}

\begin{tcolorbox}[breakable,colback=teal!4,colframe=teal!35!black,title={Before you move on}]
What does \ja{ちゅうしゃ} mean? (\textquotedblleft{}An injection\textquotedblright{}.) Say it once more.
\end{tcolorbox}
\section[\texorpdfstring{\ja{はいしゃ} (haisha)}{はいしゃ (haisha)}]{\ja{はいしゃ} (haisha) — a dentist}

\label{lesson:JA-C217-haisha}



\begin{quote}
Before the new one: say the Japanese for a consultation, then the Japanese for an injection.
\end{quote}

\subsection*{You'll want to know: \ja{はいしゃ}}
\textbf{\ja{はいしゃ}} — \emph{haisha} — \textquotedblleft{}a dentist\textquotedblright{}.

\begin{quote}
\textbf{\ja{あした} \ja{はいしゃに} \ja{いきますか。}} — \emph{ashita haisha ni ikimasu ka} — \textquotedblleft{}Are you going to the dentist tomorrow?\textquotedblright{}
\end{quote}

\begin{cousinweb}
\textbf{\ja{はいしゃ}} (\emph{haisha}) is built from native \emph{ha}, \textquotedblleft{}tooth\textquotedblright{}, inherited, and Sino-Japanese \emph{isha}, \textquotedblleft{}doctor\textquotedblright{}, borrowed from Chinese (JMdict, KANJIDIC2).
\end{cousinweb}

\begin{grammarlens}[title={What you've built}]
One more word for everyday life at A2.
\end{grammarlens}

\subsection*{Your turn}
\begin{itemize}
\raggedright
  \item \emph{Say it:} \emph{haisha}
  \item \emph{Say it:} \emph{haisha}, once more
  \item \emph{Say it:} \emph{ashita haisha ni ikimasu ka} — Are you going to the dentist tomorrow?
  \item \emph{Recall:} say the Japanese for a consultation, then the Japanese for an injection, then say \emph{haisha} again
\end{itemize}

\begin{tcolorbox}[breakable,colback=teal!4,colframe=teal!35!black,title={Before you move on}]
What does \ja{はいしゃ} mean? (\textquotedblleft{}A dentist\textquotedblright{}.) Say it once more.
\end{tcolorbox}
\section[\texorpdfstring{\ja{むしば} (mushiba)}{むしば (mushiba)}]{\ja{むしば} (mushiba) — a cavity; a decayed tooth}

\label{lesson:JA-C217-mushiba}



\begin{quote}
Before the new one: say the Japanese for an injection, then the Japanese for a dentist.
\end{quote}

\subsection*{You'll want to know: \ja{むしば}}
\textbf{\ja{むしば}} — \emph{mushiba} — \textquotedblleft{}a cavity; a decayed tooth\textquotedblright{}.

\begin{quote}
\textbf{\ja{むしばは} \ja{ありません。}} — \emph{mushiba wa arimasen} — \textquotedblleft{}I have no cavities.\textquotedblright{}
\end{quote}

\begin{cousinweb}
\textbf{\ja{むしば}} (\emph{mushiba}) is native Japanese (wago), built from \emph{mushi}, \textquotedblleft{}insect, worm\textquotedblright{}, and \emph{ha}, \textquotedblleft{}tooth\textquotedblright{}, voiced to \emph{ba}, both inherited rather than borrowed (JMdict, KANJIDIC2).
\end{cousinweb}

\begin{grammarlens}[title={What you've built}]
One more word for everyday life at A2.
\end{grammarlens}

\subsection*{Your turn}
\begin{itemize}
\raggedright
  \item \emph{Say it:} \emph{mushiba}
  \item \emph{Say it:} \emph{mushiba}, once more
  \item \emph{Say it:} \emph{mushiba wa arimasen} — I have no cavities.
  \item \emph{Recall:} say the Japanese for an injection, then the Japanese for a dentist, then say \emph{mushiba} again
\end{itemize}

\begin{tcolorbox}[breakable,colback=teal!4,colframe=teal!35!black,title={Before you move on}]
What does \ja{むしば} mean? (\textquotedblleft{}A cavity; a decayed tooth\textquotedblright{}.) Say it once more.
\end{tcolorbox}
